\lysh{34820}{2}
\begin{alist}
\item Οι πόντοι που πέτυχε ο καλαθοσφαιριστής στους τελευταίους πέντε ($ν = 5$) αγώνες της ομάδας του σε αύξουσα σειρά είναι:
\[
    17, 19, 20, 21, 23
\]
Η μέση τιμή των πόντων του καλαθοσφαιριστή (πόντοι ανά αγώνα) είναι:
\[
    \bar{x} = \frac{\sum_{i=1}^{5} t_i}{5} = \frac{100}{5} = 20
\]

\item Η διακύμανση του δείγματος υπολογίζεται από τον τύπο:
\[
    s^2 = \frac{1}{ν} \sum_{i=1}^{5} (t_i - \bar{x})^2 = \frac{1}{5} \cdot 20 = 4
\]

\item Ο συντελεστής μεταβολής του δείγματος είναι:
\[
    CV = \frac{s}{|\bar{x}|}
\]
όπου $s = \sqrt{s^2} = \sqrt{4} = 2$. Επομένως:
\[
    CV = \frac{s}{|\bar{x}|} = \frac{2}{20} = 0,10 = 10\%
\]
Αφού ο συντελεστής μεταβολής δεν ξεπερνά το 10\%, συμπεραίνουμε ότι το δείγμα είναι ομοιογενές.
\end{alist}
\vspace{6mm}
\noindent\rule{\textwidth}{0.4pt}
\vspace{6mm}

